% 2010 年湖南卷文科第 20 题：数表序列中的表 1、表 2、表 3。
% 来源 PDF：普通高考/2010/2010湖南文.pdf，SHA-256 38611c33c9da2139ad6cd4117513ee144b09a6860180eb976c446b4858ada122。
% 原页：第 2 页（A3 三栏）第 2 栏上部，书页码 604。
% 重绘依据：照原页无框格排版，表 1、表 2、表 3、⋯ 自左向右并列；
%   表 1 为 1；表 2 为「1 3」与「4」；表 3 为「1 3 5」「4 8」「12」，各行左对齐。
\begin{tikzpicture}[scale=0.95, semithick, every node/.style={inner sep=1.5pt}]

  % 表头
  \node at (0,0) {表 1};
  \node at (2.4,0) {表 2};
  \node at (5.4,0) {表 3};
  \node at (7.7,0) {$\cdots$};

  % 第 1 行
  \node at (0,-1.0) {1};
  \node at (2.0,-1.0) {1};
  \node at (3.0,-1.0) {3};
  \node at (4.7,-1.0) {1};
  \node at (5.5,-1.0) {3};
  \node at (6.3,-1.0) {5};

  % 第 2 行
  \node at (2.0,-1.8) {4};
  \node at (4.7,-1.8) {4};
  \node at (5.5,-1.8) {8};

  % 第 3 行
  \node at (4.7,-2.6) {12};

\end{tikzpicture}
